\subsection{Approximating the predictive entropy}\label{sec:predictiveEntropy}

We now show how to approximate $\H[p(y|\D_n,\x,\xopt)]$ in
(\ref{eq:newObjective}). Note that we can write the argument to $\H$ in this expression as
$p(y|\D_n,\x,\xopt) = \int p(y|f(\x)) p(f(\x)|\D_n,\xopt)\,df(\x)$.  
Here $p(f(\x)|\D_n,\xopt)$ is the posterior distribution on $f(\x)$ given 
$\D_n$ and the location $\xopt$ of the global maximizer of $f$.
When the likelihood $p(y|f(\x))$ is Gaussian, we have that
$p(f(\x)|\D_n)$ is analytically tractable since it is the predictive distribution
of a Gaussian process.  However, by further conditioning on the location $\xopt$
of the global maximizer we are introducing additional constraints, namely that $f(\vz)\leq f(\xopt)$ for all $\vz\in\X$.
These constraints make $p(f(\x)|\D_n,\xopt)$ intractable.  To circumvent this
difficulty, we instead use the following simplified constraints:
\begin{enumerate}
    \item[C1.] \textbf{$\xopt$ is a local maximum.} This is achieved by letting $\nabla
    f(\xopt)=\zero$ and ensuring that $\nabla^{2}f(\xopt)$ is negative
    definite. We further assume that the non-diagonal elements of
    $\nabla^{2}f(\xopt)$, denoted by $\upper[\nabla^{2}f(\xopt)]$, are known,
    for example they could all be zero. This simplifies the negative-definite constraint.
    We denote by C1.1 the constraint given by $\nabla f(\xopt)=\zero$ and $\upper[\nabla^2 f(\xopt)] = \zero$.
    We denote by C1.2 the constraint that forces the elements of $\diag[\nabla^2 f(\xopt)]$ to be negative.

    \item[C2.] \textbf{$f(\xopt)$ is larger than past observations.} 
    We also assume that $f(\xopt)\geq f(\x_i)$ for all $i\leq n$.
    However, we only observe $f(\x_i)$ noisily via $y_i$.
    To avoid making inference on these latent function values, we approximate the above hard constraints
    with the soft constraint $f(\xopt) > y_\mathrm{max} + \epsilon$, where
    $\epsilon\sim\Normal(0,\sigma^2)$ and $y_\mathrm{max}$ is the largest $y_i$ seen so far.
    
    \item[C3.] \textbf{$f(\x)$ is smaller than $f(\xopt)$.} This simplified
    constraint only conditions on the given $\x$ rather than requiring
    $f(\xopt)\leq f(\vz)$ for all $\vz\in\X$.
\end{enumerate}
We incorporate these simplified constraints into $p(f(\x)|\D_n)$ to 
approximate $p(f(\x)|\D_n,\xopt)$.  This is achieved by multiplying
$p(f(\x)|\D_n)$ with specific factors that encode the above constraints. In what
follows we briefly show how to construct these factors; more detail is given in
Appendix~\ref{sec:predictiveEntropyAppendix}.

\newcommand\C[1]{\mathrm{C#1}}

% MATT: we don't really need to state the following bits here as we don't make
% use of these facts.

%Let the prior distribution for $f$ be a zero-mean Gaussian process with
%covariance function given by $k[\x_{i},\x_{j}] = \gamma \exp\left\{
%-0.5(\x_{i}-\x_{j})^{\text{T}}\vL(\x_{i}-\x_{j})\right\}$, where $\vL$ is a
%diagonal matrix with the lengthscale of the covariance function.

%These variables are concatenated in the observed vector $\vo$ of dimension $n +
%d + d(d-1)/2$.

%the latent vector $\vl$ of dimension $1 + d$. 

Consider the latent variable 
    $\vz = [f(\xopt);\, \diag[\nabla^2f(\xopt)]]$.
To incorporate constraint C1.1 we can condition on the data and on the
``observations'' given by the constraints $\nabla f(\xopt)=\zero$ and
$\upper[\nabla^2 f(\xopt)]=\zero$.  Since $f$ is distributed according to a GP,
the joint distribution between $\vz$ and these observations is multivariate
Gaussian. The covariance between the noisy observations $\vy_n$ and the extra
noise-free derivative observations can easily be computed \cite{Solak:2003}. The
resulting conditional distribution is also multivariate Gaussian with mean
$\vm_0$ and covariance $\vV_0$. These computations are similar to those
performed in (\ref{eq:gpvar}). Constraints C1.2 and C2 can then be incorporated
by writing
\begin{equation}
    \textstyle
    p(\vz|\D_n,\C1,\C2) \propto
    \Phi_{\sigma^2}(f(\xopt) - y_\text{max})
    \Big[ \prod_{i=1}^{d} \I\big([\nabla^2 f(\xopt)]_{ii} \leq 0\big) \Big]
    \Normal(\vz|\vm_0,\vV_0)\,,\label{eq:conditional}
\end{equation}
where $\Phi_{\sigma^2}$ is the cdf of a zero-mean Gaussian distribution with
variance $\sigma^2$.  The first new factor in this expression guarantees that
$f(\xopt) > y_\text{max} + \epsilon$, where we have marginalized $\epsilon$ out,
and the second set of factors guarantees that the entries in $\diag[\nabla^2
f(\xopt)]$ are negative.

Later integrals that make use of $p(\vz|\D_n,\C1,\C2)$, however, will not admit
a closed-form expression. As a result we compute a Gaussian approximation
$q(\vz)$ to this distribution using Expectation Propagation (EP)
\cite{Minka:2001}.  The resulting algorithm is similar to the implementation of
EP for binary classification with Gaussian processes \cite{Rasmussen:2006}. EP
approximates each non-Gaussian factor in (\ref{eq:conditional}) with a Gaussian
factor whose mean and variance are $\widetilde{m}_i$ and $\widetilde{v}_i$,
respectively. The EP approximation can then be written as
    $q(\vz)
    \propto [ \prod_{i=1}^{d+1} \Normal(z_i|\widetilde m_i, \widetilde v_i)
    ] \Normal(\vz|\vm_0,\vV_0)$.
Note that these computations have so far not depended on $\x$, so we can compute
$\{\vm_0,\vV_0, \widetilde\vm, \widetilde\vv\}$ once and store them for later
use, where $\widetilde\vm = (\tilde{m}_1,\ldots,\tilde{m}_{d+1})$ and
$\widetilde\vv = (\tilde{v}_1,\ldots,\tilde{v}_{d+1})$.

We will now describe how to compute the predictive variance of some latent
function value $f(\x)$ given these constraints. Let $\vf=[f(\vx); f(\xopt)]$ be a
vector given by the concatenation of the values of the latent function at $\x$ and $\xopt$.
The joint distribution between $\vf$,  $\vz$, the evaluations $\vy_n$ collected so far and
the derivative ``observations'' $\nabla f(\xopt)=\zero$ and
$\upper[\nabla^2 f(\xopt)]=\zero$ is multivariate Gaussian.
Using $q(\vz)$, we then obtain the following approximation:
\begin{equation}
    \textstyle
    p(\vf|\D_n,\C1,\C2) \approx
    \int p(\vf|\vz,\D_n,\C{1.1}) \,q(\vz)\,d \vz =
    \Normal(\vf|\vm_\vf,\vV_\vf)\,.
\end{equation}
Implicitly we are assuming above that $\vf$ depends on our observations and
constraint C1.1, but is independent of C1.2 and C2 given $\vz$.  The
computations necessary to obtain $\vm_\vf$ and $\vV_\vf$ are similar to those used
above and in (\ref{eq:gpvar}).
The required quantities are similar to the ones used by EP to make predictions in the
Gaussian process binary classifier~\cite{Rasmussen:2006}.  We can then
incorporate C3 by multiplying $\Normal(\vf|\vm_\vf,\vV_\vf)$ with a factor that
guarantees $f(\x) < f(\xopt)$.  The predictive distribution for $f(\x)$ given
$\D_n$ and all the constraints can be approximated as
\begin{align}
    \textstyle
    p(f(\x)|\D_n,\C1,\C2,\C3) 
    \approx
    Z^{-1} \int \I( f_1 < f_2) \,\Normal(\vf|\vm_\vf,\vV_\vf) \,df_2\,,
    \label{eq:finalPredictive}
\end{align}
where $Z$ is a normalization constant.  The variance of the right hand size of
(\ref{eq:finalPredictive}) is given by
\begin{equation}
    v_n(\x|\xopt)
    = [\vV_\vf]_{1,1} -
    s^{-1}\beta(\beta + \alpha)
    \{  [\vV_\vf]_{1,1} - [\vV_\vf]_{1,2} \}^2\,,
\label{eq:predictiveVariance}
\end{equation}
where $s = [-1, 1]\T \vV_\vf [-1, 1]$, $\alpha = \mu / \sqrt{s}$, $\mu = [-1, 1]\T
\vm_\vf$, $\beta = \phi(\alpha) / \Phi(\alpha)$, and $\phi(\cdot)$ and $\Phi(\cdot)$
are the standard Gaussian density function and cdf, respectively.  By further
approximating (\ref{eq:finalPredictive}) by a Gaussian distribution with the
same mean and variance we can write the entropy as
$\H[p(y|\D_n,\x,\xopt)] \approx 0.5 \log [2 \pi e (v_n(\x|\xopt) +
\sigma^2)  ]$.

The computation of (\ref{eq:predictiveVariance}) can be numerically unstable
when $s$ is very close to zero.  This occurs when $[\vV_\vf]_{1,1}$ is very similar
to $[\vV_\vf]_{1,2}$.  To avoid these numerical problems, we multiply $[\vV_\vf]_{1,2}$
by the largest $0\leq \kappa \leq 1$ that guarantees that $s > 10^{-10}$. This
can be understood as slightly reducing the amount of dependence between $f(\x)$
and $f(\xopt)$ when $\x$ is very close to $\xopt$.  Finally, fixing
$\upper[\nabla^2 f(\xopt)]$ to be zero can also produce poor predictions when
the actual $f$ does not satisfy this constraint.  To avoid this, we instead fix this quantity to
$\upper[\nabla^2 f\i(\xopt)]$, where $f\i$ is the $i$th
sample function optimized in Section \ref{sec:sampling} to sample $\xopt\i$.

